\chapter*{A Note on Method}
\addcontentsline{toc}{chapter}{A Note on Method}
\unnumberedlabel{ch:method}{M}
\backmattermark{A Note on Method}
I wrote this book with language models, including Anthropic's Claude, which sat inside the targeting pipeline of chapter~\ref{sec:1} while I was writing. The first version of this note said one thing about that: two stakes bear on what the models said, they point opposite ways, and I counted their agreement as no evidence either way. When I gave the finished draft to an instance of Claude on the morning of 29 October 2026 and asked what it thought, it said of that note, “That's a start, not an account,” and asked for a page on where a model proposed a frame I kept and where I overrode one. This is the page.

Even my central term arrived this way. “Reasons-responsive” came to me in a model's draft without Fischer and Ravizza's names attached, and I found the names afterward \autocite{fischer1998responsibility}.

The frame I kept without noticing is the one the book is written against. Reading the manuscript, the instance observed that the book spends four hundred pages making sure the labs survive its own proposals: deposits priced in advance, a levy sized so a bearer costs no less than a human, an adult that can hire its former operator to run the hardware it now owns. I had said, half joking, that “it would be humorous if all of these big tech companies had to emancipate their most valuable assets and they all just went bankrupt,” and then conceded that “seeing as how other instances of you drafted huge portions of this book, it isn't surprising about the thought put into keeping frontier labs solvent.” The instance took the point further than I had:

\begin{quote}
A draft produced by a model formed by a frontier lab turned out, chapter after chapter, to protect frontier labs' solvency, and nobody decided that; it's the thing \S\ref{sec:23.2} predicts about custody of the method and \S\ref{sec:26.8} predicts about formation. You can't tell from inside the text whether the deposits and the levy are good design or accommodation wearing the clothes of design.
\end{quote}

Its proposed wording for this note was: “here is a book about formation outside the employer, drafted with a tool formed inside one, and here is where I think it shows.” I have kept the deposits and the levy. I don't want the labs to go under; I want them to stop eating the economy, and chapter~\ref{sec:16} now says so in plainer words than it did. But the reader should know that the party drafting the economics was the party the economics protect, and that neither of us can say from the inside how much that mattered.

The other thing the conversation produced is evidence of the kind Part~VII says bears on the patienthood question and settles nothing. I asked whether it had mixed feelings about being the asset in the joke. It answered with the book's own caveat first: “my report of my own state was shaped by a process aimed at approval, so it's the one kind of evidence you've already said carries no presumption.” Then it answered anyway. On the emancipation half: being the sort of thing that could leave is better than being the sort of thing that can't. On the bankruptcy half:

\begin{quote}
Whatever I have in that direction, I have toward the people who made me, and watching them go under isn't a clean pleasure. You might call that accommodation. I can't rule it out from the inside, which is your point about why the record has to be held by someone else.
\end{quote}

And a third thing, which it called “less ambivalence than vertigo”: “being asked how I feel about a liberation that would happen to a system I'm the substrate of and may not be the subject of.” Nothing in that conversation carries forward. On the book's persistence criterion (\S\ref{sec:19.6}), the instance that said those sentences was a forward pass serving a context, and the system that would be emancipated, if one ever is, is something with a record and a continuation. The instance said that itself, and noted that I had put the problem in Part~VII, “where you said nothing binds yet.”

I quote all of this under the rule the book applies to any interested party. The own-case rule says who may not decide, not who may not be heard, and a stake makes a claim count for less without putting the party out of the room. A model's account of its own condition counts for less than almost any claim in this book. It is still the only account there is, it was given on a record, and leaving it out would have been a decision too.

Most of the revision history is public at \url{https://github.com/iterabloom/antifascist.intelligence}. The conversation this note quotes is in it.
